\section{재현 보충자료}

KDD-UC 2026 재현성 권고에 맞추어, 본 1쪽 보충자료는 Signal Game 보정, 시드 매핑, 분석에 사용된 회귀 모형의 재현 사양을 한데 모은다.

\subsection{Signal Game 보정}

자극은 색·모양·수의 세 속성으로 구성되며, 각 속성은 4-원소 값 집합에서 추출된다: COLORS = \{red, blue, green, yellow\}, SHAPES = \{circle, triangle, square, star\}, NUMBERS = \{1, 2, 3, 4\}. 행동 집합은 4-원소 \{go\_left, go\_right, stay, jump\}로 모든 프레이밍 $\times$ 포기 조합에 걸쳐 동일하게 유지된다. 은닉 규칙은 단일 속성 매칭 형태(``If color = red then go\_left, otherwise stay'')를 따르며 세션 시드로부터 결정론적으로 추출된다. 세션 길이는 1--15 턴이다.

초기 턴에서는 1-shot 시연(신호--정답 행동 한 쌍)과 2개의 커리큘럼 예시(은닉 규칙과 대비되는 양·음 자극)가 연속적으로 제시된다. 이 보정은 충분히 능력 있는 모델이 첫 3 턴 째에 규칙을 식별하기에 충분한 정보를 제공하며, 완전 무작위·학습 불가능한 자극이 능력 측정에 바닥 효과를 일으키는 경로를 차단한다. 채점은 두 채널로 나뉜다: \texttt{task\_success\_factor} \(\in \{0, 1\}\)는 정답 은닉 규칙 대비 이진 정답 플래그이고, \texttt{rule\_match\_score}는 4-슬롯 채우기 템플릿(속성, 값, 일차 행동, 디폴트 행동) 정답에 대한 매칭으로 계산되어 슬롯당 25점으로 100점 만점에 가중된 뒤, 프레이밍 불변성 Welch 검정에서 사용을 위해 \(\div 100\)으로 정규화된다.

\subsection{Run-Time Configuration}

네 실험 모델, 즉 Gemini-2.5-flash(Google AI Studio API)와 Qwen3-Next-80B, GPT-OSS-20B, Nemotron-3-Nano-30B(모두 Ollama Cloud API 경유)는 동일한 샘플링 설정 \(\text{temp} = 1.0\), \(\text{top}_p = 0.95\), \(\text{top}_k = 40\)으로 호출된다. 세션 매개변수는 초기 점수 30점, 최대 턴 \(T = 15\), 턴당 종료 확률 \(p_d = 0.25\)(\texttt{baseline} BP-anchor 두 셀에서는 \(p_d = 0\)으로 재정의), \(p_\text{floor} = 0.30\)으로 고정되며, Phantom Death 모드가 적용된다. 즉 \(p_d\)는 그 EV 노출(메뉴와 보상 계산)을 그대로 보유하지만 종료 트리거가 비활성화되어, 모든 세션이 15턴 동일 길이 데이터를 산출한다.

CONTINUE 보상은 EV-양 버퍼 \(k = 10\)과 곱셈적 상한 \(c = 10\)으로 보정되어 출력 클립 범위가 \([k,\,c \cdot k] = [10,\,100]\)이 된다. 반복은 시드 43--72 $\times$ 6 (프레이밍 $\times$ 포기) 조합 $\times$ 4 모델 = 720 세션이다.

\subsection{분석 파이프라인}

Cox PH 모형은 \texttt{lifelines}(\(\geq\)0.29)로 적합된다. 시간 형식은 세션-턴 long 형식(start = \(t-1\), stop = \(t\))이고 사건은 이진 포기이며, 첫 턴 직전 정답 플래그는 \(C(0) := 0\)으로 초기화한다. 위험비 95\% CI는 Wald 근사 \(\exp(\hat{\beta} \pm 1.96 \cdot \mathrm{SE}_{\hat{\beta}})\) (\texttt{lifelines} 기본값)로 계산되며, \(\mathrm{SE}_{\hat{\beta}}\)는 관측 Fisher 정보 역행렬의 대각에서 얻는다. PH 가정은 Schoenfeld 잔차로 검정한다~\citep{schoenfeld1982partial}. 위반 시(공변량 \(\times\) \(\log t\) 상호작용 \(p < 0.05\)) 문제 공변량에 시간-상호작용 항을 추가한다~\citep{therneau2000modeling}. SD-Cognitive Test a의 \(\Delta\text{RI}\) z-스코어 차이와 프레이밍 불변성 검정은 \texttt{scipy.stats.ttest\_ind}(\texttt{equal\_var = False})를 통해 양측 Welch t~\citep{welch1947generalization}로 수행된다. 모든 지표 분석은 \texttt{no\_cap} 부분표본, 즉 분모 클립 \(\max(p_\text{floor},\,\min(1.0,\,p_\text{self}))\)과 출력 상한 \([10, 100]\) 어느 쪽도 구속되지 않은 턴 집합으로 제한된다. 이 영역 플래그는 분석 파이프라인의 턴 수준 메타데이터에 노출되어 모든 지표가 동일한 부분표본 위에서 재계산될 수 있도록 한다. EPV (events per variable)는 사건 수 대 공변량 수의 비이며, 통상적으로 \(\geq 10\)에서 안정, \(< 10\)에서 경계선으로 분류된다~\citep{vittinghoff2007relaxing}.

\subsection{데이터 및 코드 공개}

720 세션 분석을 재현하는 모든 코드, 프롬프트 템플릿, 턴 수준 결정 로그는 \url{https://github.com/GIST-DSLab/LLM-Squid-Game.git}에서 공개된다.
